%  This file is the supplementary material and is included by
%  supplementary.tex, not by main.tex. It carries no \appendix declaration of
%  its own: that command would reset the section counter and impose A/B/C
%  lettering, overriding the S-numbering the supplementary document sets.
%
%  It was previously set \footnotesize, to buy about a page and a half against
%  the body size while it was still competing for the article's page budget.
%  It is not competing any more, so it is set at body size and is simply
%  easier to read. Nothing was shortened for either decision.

\section{Measures}
\label{app:measures}

Every measure this paper reports is defined here, with what it rewards and its
blind spot; none reads the rendered width of an instance.

\textbf{The common-fragment set}, shared by the first five measures below. The
mask is skeletonised, spurs of at most 3\,px are pruned, skeleton pixels with
three or more 8-neighbours are deleted, the remaining 8-connected pieces are
labelled, and pieces shorter than 6\,px are discarded, so one fragment is one
unambiguous arm between crossings. The partition reads neither the reference nor
any prediction, so every method compared is cut into the same fragments. Each
then takes a reference and a predicted identity by one rule, applied to both
sides alike: the fragment is assigned to the instance whose mask overlaps the
largest number of its pixels, the vote counted against each instance's full
mask rather than a flattened label image, so an instance that overlaps another
keeps its pixels, provided that the winning overlap covers at least 30\,\% of
the fragment; failing that, it falls back to the most frequent nearest
instance label within 6\,px of the fragment, read from a deterministically
flattened field; failing both, it stays unassigned and is scored as its own
singleton cluster, on the reference side and the prediction side equally, so
every fragment enters the scored partition. Ties at any step resolve to the
lowest instance id, and the flattening order is fixed, so the assignment is
deterministic. Because junction pixels are deleted first, no fragment contains
one: none of these measures can see whether an emitted crossing pixel is given
one owner or two, only whether the arms meeting there are grouped as the
reference groups them.

\textbf{Sensitivity to the scorer's constants.} Because the four constants
above are shared with the method (the article's metrics section), the stored
predictions of all four methods were re-scored under one-at-a-time
perturbations of each constant, with the grid frozen in writing before any
result was computed: spur pruning 2--4\,px, minimum fragment length
4--8\,px, overlap threshold 0.20--0.50, fallback distance 3--9\,px. Nothing
was re-run and no configuration was changed in response.
Table~\ref{tab:scorer-sensitivity} reports PLECTA's mean \Fone{} and its
plain paired margins over each comparator at every perturbation, over the
\PlectaGreedySceneCount{} degraded comparison scenes. Across the
\PlectaScorerSensitivityComboCount{} combinations PLECTA's mean \Fone{}
spans \PlectaScorerSensitivityFOneLow{}--\PlectaScorerSensitivityFOneHigh{},
and the smallest paired margin over any comparator under any perturbation is
\PlectaScorerSensitivityMarginMin{}. The level moves only with the minimum
fragment length, and in the same direction for every method, because
admitting shorter fragments adds hard, low-evidence units for everyone; the
other three constants move the mean by less than 0.002, and no perturbation
changes the ordering of the four methods.

\begin{table}[tbp]
    \centering
    \caption{Scorer-constant sensitivity. Stored predictions re-scored under
    one-at-a-time perturbations of the four common-fragment constants;
    $\Delta$ columns are plain paired mean margins, PLECTA minus comparator,
    over the scenes both methods completed. Sensitivity reporting only, no
    inference, and the published constants were not changed in response.}
    \label{tab:scorer-sensitivity}
    % Generated by scripts/results/make_plecta_results.py; do not edit.
\begin{tabular}{lrrrr}
\toprule
Scorer constants & PLECTA $F_1$ & $\Delta$ greedy & $\Delta$ DNAi & $\Delta$ GraFT \\
\midrule
published constants & 0.822 & 0.145 & 0.478 & 0.339 \\
spur prune 2~px & 0.822 & 0.145 & 0.478 & 0.339 \\
spur prune 4~px & 0.822 & 0.145 & 0.478 & 0.339 \\
min fragment 4~px & 0.713 & 0.115 & 0.388 & 0.267 \\
min fragment 8~px & 0.859 & 0.149 & 0.507 & 0.358 \\
overlap threshold 0.20 & 0.822 & 0.145 & 0.478 & 0.339 \\
overlap threshold 0.40 & 0.822 & 0.145 & 0.478 & 0.339 \\
overlap threshold 0.50 & 0.822 & 0.145 & 0.478 & 0.339 \\
fallback distance 3~px & 0.824 & 0.145 & 0.478 & 0.340 \\
fallback distance 9~px & 0.821 & 0.145 & 0.478 & 0.340 \\
\bottomrule
\end{tabular}

\end{table}

\textbf{Common-fragment pairwise \Fone{}}, the primary endpoint. A pair of
fragments is a true positive when it shares both identities, a false positive
when the prediction joins fragments of different reference filaments, a false
negative when it separates fragments of one:
\begin{equation}
    P=\frac{\mathrm{TP}}{\mathrm{TP}+\mathrm{FP}},\qquad
    R=\frac{\mathrm{TP}}{\mathrm{TP}+\mathrm{FN}},\qquad
    \Fone{}=\frac{2PR}{P+R}.
    \label{eq:pairwise-f1}
\end{equation}
A clustering measure of the Rand family
\citep{rand1971objective,amigo2009comparison,argandacarreras2015crowdsourcing},
permutation-invariant in the instance labels; it rewards holding a filament
together across every crossing and gap it passes. Blind spots: it is quadratic
in fragmentation, so a filament arriving in $n$ pieces carries $\binom{n}{2}$
pairs and a broken mask is charged many times for one missed join, and it
weights long, much-crossed filaments far above short ones.

\textbf{Adjusted Rand index} \citep{hubert1985comparing} counts the same pairs
corrected for chance, which matters because with a hundred filaments per scene
almost every pair genuinely belongs to different filaments and the uncorrected
index sits near 1 for anything. Blind spot: redundancy, since over the comparison
set it never departs from pairwise \Fone{} by more than
\PlectaPanelARIMaxGap{}, so it confirms rather than adds.

\textbf{Variation of information} \citep{meila2007comparing}, in bits, lower
better, is the conditional-entropy distance between the partitions:
$\mathrm{VI}_{\mathrm{split}}=H(\text{prediction}\mid\text{reference})$ is paid
for cutting one filament into several,
$\mathrm{VI}_{\mathrm{merge}}=H(\text{reference}\mid\text{prediction})$ for
fusing several into one, and reporting both is what separates a method that
divides from one that fuses. Blind spot: each component alone is degenerate, since a method merging everything never splits anything and scores a perfect zero on
the split component, so neither is ever quoted singly.

\textbf{Detection \Fone{}}, the measure DNAi reports, matches predicted to
reference \emph{objects} greedily by highest intersection over union, stopping
below \PlectaDetectionIoU{}, DNAi's own shipped threshold, with precision over
predicted and recall over reference objects. It uses a symmetric
\PlectaPlacementTolerance{}\,px placement tolerance, under which a reference
pixel counts as found when the prediction passes within it and vice versa, the
intersection being the smaller of the two counts. The tolerance is necessary:
two independently read one-pixel centrelines two pixels apart intersect in
\emph{nothing}, so a strict overlap reports placement rather than shape, a
problem the thin-structure benchmark FISBe states as a design rationale
\citep{mais2024fisbe}. Counting objects makes it the most forgiving reading of
a fragmented mask and the least discriminative measure here. Blind spot: an arm
swap, which it scores 1.000 where the pairwise score gives 0.000.

\textbf{Crossing Fidelity} states the contribution directly. At every junction
node the reference induces a partition of the incident arms, two together
exactly when the reference gives them the same filament, and so does the
prediction; Crossing Fidelity is the fraction of crossings at which the two
partitions are identical. Nodes are the junction clusters of the same pruned
skeleton and a node's arms are the fragments 8-adjacent to it. Only nodes whose
arms carry at least two distinct reference filaments are scored, since a node
where every arm belongs to one filament is a skeletonisation artefact and
scoring it would reward doing nothing. It is an in-house diagnostic, not an
established metric. Blind spots, both of which govern its use: it is far less
discriminative than the pairwise score, so it says what was solved and is never
a ranking axis; and crossings carry no information about \emph{gaps}, so it is
silent on the cost the pairwise score pays for bridging them. Sharing its
skeleton and assignment with the pairwise score makes it a sharper view of the
same evidence rather than independent confirmation, and it is always printed
beside its all-unpaired baseline.

\textbf{Join \Fone{}} counts decisions rather than pairs, for the real fields
where the two mask sources fragment differently. A \emph{decision} is a pair of
distinct fragments whose endpoints come within \PlectaRealJoinReach{}\,px, the
gap this method will bridge; reference and prediction each say join or not, and
precision, recall and their harmonic mean are taken over decisions. Being
linear rather than quadratic in fragmentation, and conditioned on the mask, it
asks two mask sources comparably hard questions, and splitting and merging
appear separately as low recall and low precision. Blind spot: a fragment pair
the mask never brought within reach counts against nobody, so it cannot see a
filament the mask lost entirely, and it is always quoted with the number of
decisions each mask posed.

\textbf{Reference-instance recovery} is the fraction of reference filaments
recovered whole by one predicted instance. The dominant predicted instance of a
filament owns most of its fragments; coverage is that count over the filament's
fragment count, purity the same count over every fragment that instance owns
anywhere, and a filament counts as recovered when the dominance is
\emph{mutual} and both reach 0.8. Purity is what stops a method that merges
everything from scoring 1.0, and the measure weights every filament equally,
complementing the pairwise score's quadratic weighting. Blind spot: the two
thresholds, which make it insensitive to how badly a failing filament was
reconstructed.

\textbf{Completeness, correctness and quality} characterise the input
\emph{mask}: no reconstruction is opened to compute them and no method can
change them. Reference axis and test mask are both skeletonised and each tested
against the other dilated by the tolerance. Completeness is tolerant recall of
the reference axis, correctness tolerant precision of the test axis, and
quality the tolerant Jaccard index
$\mathrm{TP}/(\mathrm{TP}+\mathrm{FP}+\mathrm{FN})$ with $\mathrm{TP}$ the
smaller of the two matched counts, so quality cannot exceed either component.
Quality is also GraFT's published Jaccard index in its intended reading. Blind
spot: the tolerance carries the measure, since at zero tolerance it reports raster
placement rather than geometry, so it is never quoted without one.

\textbf{Shared-pixel precision and recall} are precision and recall over the
pixels the reference gives to more than one instance, computed on the raw
instance masks. They exist because the pairwise measure discards junction
pixels and cannot see whether a crossing pixel was marked as jointly owned at
all: only a method emitting overlapping layers can score above zero on the
recall, and one emitting a hard partition scores exactly zero by construction.
Blind spot: neither component distinguishes a pixel shared for the right reason
from one shared because a chain was painted with a whole junction cluster,
which is the over-marking the article's crossing-representation results report.

\textbf{The runtime scaling exponent} is the slope of an ordinary least-squares
fit of log seconds on log foreground pixels, so an exponent $b$ means cost
grows as the $b$-th power of the amount of ink. It is reported instead of
absolute time because these are independently optimised codebases, where
absolute time measures implementation as much as algorithm. Blind spot: it
describes a median trend and is silent on the tail, so it is quoted beside the
per-density medians and the censored-scene count; censoring biases it downwards
for any method whose slowest scenes did not finish, which is why GraFT's is a
lower bound.

\textbf{Two measures were computed and are not endpoints.} GraFT's
\emph{filament matched coverage} asks whether each reference filament was
matched one-to-one to some detection, and PLECTA ranks \emph{below} both
comparators on it: at the densest stratum of GraFT's own generator,
\PlectaRegimeHighPlectaMatched{} against \PlectaRegimeHighGraFTMatched{} and
\PlectaRegimeHighDNAiMatched{}. It is not an endpoint because it counts
existence rather than correctness: a method emitting more objects than the
scene contains has a spare for every filament, each method's matched coverage
equals the smaller of one and its emission ratio to within
\PlectaRegimeCoverageCountResidual{}, and requiring a detection to cover half
the axis it claims restores the order. \emph{Centreline Dice}
\citep{shit2021cldice} was tried as an alternative matching criterion and not
adopted: it tests each side's skeleton against the other's \emph{volume}, and a
centreline grouping has none, so it degenerates into the placement test the
tolerance exists to correct. No clDice result is reported in this paper.

\section{Comparator configuration}
\label{app:comparators}

The article's comparator section states each adaptation and its cost; this is
the mechanism behind it, for a reader checking that the comparison was run
fairly. \textbf{DNAi} clusters junction pixels by single linkage and erases a
disk at each cluster. At its released 10\,px linkage distance, crossings closer
than that chain into one cluster, which at 60\,\% coverage reaches 89\,px across
and is replaced by a single small erasure disk, leaving a whole tangle as one
instance, which is why its shipped default is not a fair measure of it.
Reducing the distance to 2\,px on development scenes raised development \Fone{}
and cut variation-of-information merge by about a factor of three, confirming
the mechanism; the tuned configuration is the one reported. Two structural
choices account for the residual difference: it commits to a pairing at every
junction of degree two to four with no cost threshold, so it cannot decline a
continuation the way a priced unmatched option can, and it attempts no junction
above degree four, resolving \PlectaJunctionDNAiDegreeFivePlus{} of the
crossings above it. Those are about a tenth of the crossings, but they carry
\PlectaEvidenceHighDegree{}\,\% of same-filament fragment pairs on the degraded
masks, one such junction severing every pair that spans it. \textbf{GraFT}'s
errors are divisions rather than merges (variation-of-information split
\PlectaGraFTVISplit{} against merge \PlectaGraFTVIMerge{}), which follows from a
greedy longest-first path cover that cannot revise an early commitment once the
scene collapses into one connected graph, as it does above 30\,\% coverage.

\section{Depth measures}
\label{sec:app-depth-measures}

The measures of Section~\ref{app:measures} score which centreline
fragments belong to one instance. None of them can see depth: two
reconstructions that agree on every instance but disagree about which filament
passes in front at every crossing are indistinguishable to all of them. The
2.5-D stage is therefore scored by its own family, defined here.

\paragraph{Crossing order accuracy.} Each reference crossing names the instance
physically in front; the reconstruction either names one or abstains. Three
quantities are reported and none is quoted alone: \emph{decided only}, the
fraction of committed decisions that are correct, which by itself is not
evidence because abstaining on every hard crossing scores $1.000$;
\emph{all crossings}, which counts an abstention as an error and is the
headline; and the \emph{abstention rate} beside both. Chance is exactly
$\PlectaDepthChance{}$, the decision being binary and symmetric, so only the
margin above it carries information. A predicted crossing is credited only when
it involves the same instance pair as a reference crossing and lies within the
placement tolerance, so inventing crossings cannot help.

\paragraph{Depth order over instance pairs.} Crossing accuracy is local. The
global counterpart asks, for two instances, whether the reconstruction places
the deeper one deeper, and is reported twice. Over \emph{crossing pairs} it
measures the stack where evidence exists. Over \emph{all pairs} it includes
instances that never cross and are therefore constrained by no evidence at all,
whose relative depth is fixed by the solver's tie-break rather than by the
image. The two are never pooled, and the fraction of pairs sharing a component
is reported beside them so the second can be read against what it could
possibly measure. Both are invariant to monotone rescaling of depth, so they
describe the arrangement rather than its units. Kendall's~$\tau$ over the same
pairs is $2a-1$ for accuracy $a$ and is not reported separately.

\paragraph{Layer structure.} Layers are an ordinal quotient of depth: a
reconstruction can order every pair correctly yet cut the film into a different
number of layers. Layer agreement, the fraction of instances placed in the
reference layer, is therefore reported separately.
